
\documentclass{article}
\usepackage{colortbl}
\usepackage{makecell}
\usepackage{multirow}
\usepackage{supertabular}

\begin{document}

\newcounter{utterance}

\centering \large Interaction Transcript for game `cladder', experiment `full\_v1.5\_default', episode 8179 with refineonly{-}Qwen3.5{-}27B.
\vspace{24pt}

{ \footnotesize  \setcounter{utterance}{1}
\setlength{\tabcolsep}{0pt}
\begin{supertabular}{c@{$\;$}|p{.15\linewidth}@{}p{.15\linewidth}p{.15\linewidth}p{.15\linewidth}p{.15\linewidth}p{.15\linewidth}}
    \# & \multicolumn{2}{c}{Player} && \multicolumn{2}{c}{Game Master} \\
    \hline

    \theutterance \stepcounter{utterance}  
    & & & \multicolumn{4}{p{0.6\linewidth}}{
        \cellcolor[rgb]{0.9,0.9,0.9}{
            \makecell[{{p{\linewidth}}}]{
                \texttt{\tiny{[P1$\langle$GM]}}
                \texttt{You are an expert in causal inference. The following question is not a typical commonsense query, but rather a meticulously designed question created by a professor specializing in causal inference, intended to assess the students' mastery of the course content.} \\
\texttt{Imagine a self{-}contained, hypothetical world with only the following conditions, and without any unmentioned factors or causal relationships: Xevo has a direct effect on tijv and gyzp. Tijw has a direct effect on tijv and gyzp. Tijv has a direct effect on gyzp. Tijw is unobserved. Method 1: We look directly at how xevo correlates with gyzp in general. Method 2: We look at this correlation case by case according to tijv. To understand how xevo affects gyzp, is it more correct to use the Method 1 than Method 2?} \\
\texttt{Start your answer with 'Yes' or 'No', followed by additional reasoning or evidence to support your explanation.} \\
            }
        }
    }
     \\ \\

    \theutterance \stepcounter{utterance}  
    & & \multicolumn{4}{p{0.6\linewidth}}{
        \cellcolor[rgb]{0.95,0.95,0.95}{
            \makecell[{{p{\linewidth}}}]{
                \texttt{\tiny{[P1$\leftrightarrow$P]}}
                \texttt{{[}DRAFT{]}} \\
\texttt{Yes} \\
\\ 
\texttt{**Reasoning:**} \\
\\ 
\texttt{To determine the correct approach, we must analyze the causal structure (DAG) and the role of the variable \$tijv\$ in the relationship between \$xevo\$ and \$gyzp\$.} \\
\\ 
\texttt{1.  **Causal Structure Analysis**:} \\
\texttt{    *   \$xevo \textbackslash{}rightarrow tijv\$} \\
\texttt{    *   \$xevo \textbackslash{}rightarrow gyzp\$ (Direct effect)} \\
\texttt{    *   \$tijw \textbackslash{}rightarrow tijv\$} \\
\texttt{    *   \$tijw \textbackslash{}rightarrow gyzp\$} \\
\texttt{    *   \$tijv \textbackslash{}rightarrow gyzp\$} \\
\texttt{    *   \$tijw\$ is unobserved.} \\
\\ 
\texttt{2.  **Identifying the Pathways**:} \\
\texttt{    *   **Direct Path**: \$xevo \textbackslash{}rightarrow gyzp\$.} \\
\texttt{    *   **Indirect Path**: \$xevo \textbackslash{}rightarrow tijv \textbackslash{}rightarrow gyzp\$.} \\
\texttt{    *   **Confounding Path**: \$xevo \textbackslash{}leftarrow (\textbackslash{}text\{none\}) \textbackslash{}rightarrow tijw \textbackslash{}rightarrow gyzp\$? No, \$xevo\$ and \$tijw\$ are not connected directly. However, \$tijw\$ affects both \$tijv\$ and \$gyzp\$.} \\
\texttt{    *   Crucially, \$tijv\$ is a **mediator** on the path from \$xevo\$ to \$gyzp\$. It is also a **collider** on the path \$tijw \textbackslash{}rightarrow tijv \textbackslash{}leftarrow xevo\$? No, \$xevo \textbackslash{}rightarrow tijv\$ and \$tijw \textbackslash{}rightarrow tijv\$. Yes, \$tijv\$ is a collider for the paths coming from \$xevo\$ and \$tijw\$.} \\
\\ 
\texttt{3.  **Evaluating Method 1 (Marginal Correlation)**:} \\
\texttt{    *   Method 1 estimates the **Total Effect** of \$xevo\$ on \$gyzp\$.} \\
\texttt{    *   This includes the direct effect (\$xevo \textbackslash{}rightarrow gyzp\$) and the indirect effect through the mediator (\$xevo \textbackslash{}rightarrow tijv \textbackslash{}rightarrow gyzp\$).} \\
\texttt{    *   Since \$tijw\$ is unobserved, does it create bias? \$tijw\$ affects \$tijv\$ and \$gyzp\$. However, \$xevo\$ is not caused by \$tijw\$. The path \$xevo \textbackslash{}rightarrow tijv \textbackslash{}leftarrow tijw \textbackslash{}rightarrow gyzp\$ is blocked at the collider \$tijv\$ unless we condition on it. Therefore, in the marginal distribution (Method 1), the unobserved confounder \$tijw\$ does not create a spurious association between \$xevo\$ and \$gyzp\$. The correlation observed in Method 1 represents the true total causal effect.} \\
\\ 
\texttt{4.  **Evaluating Method 2 (Conditioning on \$tijv\$)**:} \\
\texttt{    *   Method 2 involves stratifying or conditioning on \$tijv\$.} \\
            }
        }
    }
    & \\ \\

    \theutterance \stepcounter{utterance}  
    & \multicolumn{4}{p{0.6\linewidth}}{
        \cellcolor[rgb]{0.561,0.737,0.561}{
            \makecell[{{p{\linewidth}}}]{
                \texttt{\tiny{[P1$\rangle$GM]}}
                \texttt{Yes} \\
\\ 
\texttt{**Reasoning:**} \\
\\ 
\texttt{To determine the correct approach, we must analyze the causal structure (DAG) and the role of the variable \$tijv\$ in the relationship between \$xevo\$ and \$gyzp\$.} \\
\\ 
\texttt{1.  **Causal Structure Analysis**:} \\
\texttt{    *   \$xevo \textbackslash{}rightarrow tijv\$} \\
\texttt{    *   \$xevo \textbackslash{}rightarrow gyzp\$ (Direct effect)} \\
\texttt{    *   \$tijw \textbackslash{}rightarrow tijv\$} \\
\texttt{    *   \$tijw \textbackslash{}rightarrow gyzp\$} \\
\texttt{    *   \$tijv \textbackslash{}rightarrow gyzp\$} \\
\texttt{    *   \$tijw\$ is unobserved.} \\
\\ 
\texttt{2.  **Identifying the Pathways**:} \\
\texttt{    *   **Direct Path**: \$xevo \textbackslash{}rightarrow gyzp\$.} \\
\texttt{    *   **Indirect Path**: \$xevo \textbackslash{}rightarrow tijv \textbackslash{}rightarrow gyzp\$.} \\
\texttt{    *   **Backdoor Path**: There is a path \$xevo \textbackslash{}rightarrow tijv \textbackslash{}leftarrow tijw \textbackslash{}rightarrow gyzp\$. Note that \$tijv\$ acts as a **collider** on the path between \$xevo\$ and \$tijw\$.} \\
\texttt{    *   Crucially, \$tijv\$ is a **mediator** for the effect of \$xevo\$ on \$gyzp\$.} \\
\\ 
\texttt{3.  **Evaluating Method 1 (Marginal Correlation)**:} \\
\texttt{    *   Method 1 estimates the **Total Effect** of \$xevo\$ on \$gyzp\$.} \\
\texttt{    *   This includes the direct effect (\$xevo \textbackslash{}rightarrow gyzp\$) and the indirect effect through the mediator (\$xevo \textbackslash{}rightarrow tijv \textbackslash{}rightarrow gyzp\$).} \\
\texttt{    *   Regarding confounding: The unobserved variable \$tijw\$ affects both \$tijv\$ and \$gyzp\$. However, the path \$xevo \textbackslash{}rightarrow tijv \textbackslash{}leftarrow tijw \textbackslash{}rightarrow gyzp\$ is blocked at the collider \$tijv\$ because we are not conditioning on it. Since there are no other open backdoor paths from \$xevo\$ to \$gyzp\$, the marginal correlation in Method 1 provides an unbiased estimate of the total causal effect.} \\
\\ 
\texttt{4.  **Evaluating Method 2 (Conditioning on \$tijv\$)**:} \\
\texttt{    *   Method 2 involves stratifying or conditioning on \$tijv\$.} \\
\texttt{    *   **Collider Bias**: Conditioning on \$tijv\$ (a collider on the path \$xevo \textbackslash{}rightarrow tijv \textbackslash{}leftarrow tijw\$) opens a spurious path between \$xevo\$ and \$tijw\$. Since \$tijw\$ also causes \$gyzp\$, this creates a non{-}causal association between \$xevo\$ and \$gyzp\$ (specifically, \$xevo \textbackslash{}rightarrow tijv \textbackslash{}leftarrow tijw \textbackslash{}rightarrow gyzp\$ becomes open). This introduces bias} \\
            }
        }
    }
    & & \\ \\

    \theutterance \stepcounter{utterance}  
    & & & \multicolumn{2}{p{0.3\linewidth}}{
        \cellcolor[rgb]{0.9,0.9,0.9}{
            \makecell[{{p{\linewidth}}}]{
                \texttt{\tiny{[GM$|$GM]}}
                \texttt{Yes} \\
            }
        }
    }
    & & \\ \\

    \theutterance \stepcounter{utterance}  
    & & & \multicolumn{2}{p{0.3\linewidth}}{
        \cellcolor[rgb]{0.9,0.9,0.9}{
            \makecell[{{p{\linewidth}}}]{
                \texttt{\tiny{[GM$|$GM]}}
                \texttt{yes} \\
            }
        }
    }
    & & \\ \\

    \theutterance \stepcounter{utterance}  
    & & & \multicolumn{2}{p{0.3\linewidth}}{
        \cellcolor[rgb]{0.9,0.9,0.9}{
            \makecell[{{p{\linewidth}}}]{
                \texttt{\tiny{[GM$|$GM]}}
                \texttt{game\_result = WIN} \\
            }
        }
    }
    & & \\ \\

\end{supertabular}
}

\end{document}
